\chapter{CDCL}

Conflict Driven Clause Learning (CDCL), is \textit{the} modern SAT solving 
algorithm. It builds off of the ideas of DPLL and takes it to a whole new level.
The ``conflict driven'' part of CDCL means that we analyze conflicts (contradictions)
and see what we can learn from them. The ``clause learning'' part is exactly what it 
sounds like, we add new clauses to our database. In other words, CDCL is about 
analyzing why conflicts occur when we solve and using those conflicts to learn 
new clauses that will help enforce new rules as to what values variables can take,
essentially changing the CNF equation into one that's logically equivalent, but 
easier to solve. I won't lie to you, I was not able to implement this algorithm the 
first time I tried. I hope and pray that this time I will be successful, so that 
\textit{we} will be successful together dear reader. 

\textit{Aside: why have we been opting for indices so often? You see, although
we live in the land of pseudocode or machine-agnosticism, implementations don't
(duh). So in our previous chapter in DPLL, there were times where we used clause 
indices instead of pointers to clauses (obviously using copies would be terrible 
for performance and memory), but why did we use indices? Now that we're implementing 
CDCL, we're going to be adding clauses to the database quite often, and that can 
trigger a resize on the database. If we used pointers to clauses, after a resize 
triggers, the memory address of each clause would change, thus making the pointers 
invalid. We would need to add new code after the insertion of each clause to make 
sure a resize didn't trigger, and if one did, we'd need to update each pointer for 
each reason clause and watched literal, which would be very difficult because after 
a resize, we now have no way to tell which clauses mapped to which reasons and watched 
literals, because, well, the old mapping isn't valid and contains no information about 
the new addresses for everything.
}

\footnote{More information about CDCL can be found here: 
    \url{https://www.cis.upenn.edu/~cis1890/files/Lecture4.pdf}}
\section{Clause Learning}

The first step in creating a CDCL solver is to learn and use new clauses. There are 
many versions involving different ways of backtracking, variable assigning, and even 
how we learn new clauses, but for now, we will keep our implementation as close to 
our iterative DPLL version and work our way up to the more advanced and widely used 
versions of the algorithm. 

In order to learn a new clause, we first need to understand the learning process. 
The big idea is that we will encounter conflicts and contradictions while assigning 
and propagating new literals, and we want to learn from those conflicts and contradictions 
in a way that would ideally prevent us from making the same mistake again. Our unit 
propagation method is the perfect tool for this because it already analyzes our clauses 
for us and can detect when a contradiction has occurred. Rather than returning True or 
False to indicate a successful or unsuccessful propagation respectively, when we 
encounter a contradiction, we want to return the clause that is unsatisfied. You may know that
when we propagate, there could be more than one clause that is unsatisfied, but just not 
yet evaluated to be unsatisfied. Even though there could be many conflicting clauses at
once, we just want to return the first one that we find. You may wonder why not learn from 
all of them, and the answer is that by resolving the bad decision/propagation we made, 
we will (hopefully) prevent ourselves from doing more propagations that lead to those 
extra unsatisfied clauses.

After we return the conflicting clause to the main solving method, we then want to 
perform ``clause resolution.'' As stated in the beginning, there are many versions of 
clause resolution, mainly differing in when they/how they stop resolving. For our basic 
CDCL implementation, we will resolve the conflicting clause until it contains no more 
propagated literals, ie, it's a clause of only decisions. The reason why this is a good 
stopping point is because we made all decisions, and because of that, clause learning 
will essentially allow us to learn which decisions we can't make together. 

Lastly, after we resolve on the learned clause, we want to add it to the database,
and we want to initialize its watched literals so that we can use it as if it were 
part of our initial CNF equation. Before we discuss how clause resolution works at 
a theory level, let's first discuss a few implementation changes we'll need to make.

\begin{lstlisting}
    public interface:
        SATSolver(string input)
        bool solveCNF()

    private interface:
        enum var
        {
            False
            True
            Unassigned
        }

        TrailEntry {
            int lit,
            int decisionLevel,
            int clauseIdx
        }

        DecisionLevel {
            int idx,
            bool triedOtherBranch
        }

        void parse(string equation)
        bool propagate(int[] decisionLevel, int decisionLevel)
        bool createWatched()
        bool createWatched(int idx, int decisionLevel)

        int numVars
        int numClauses
        int trailHead

        int[][] clauses
        
        TrailEntry[] trail
        DecisionLevel[] trailStarts
        int[] var_to_trail

        optional var[] vars
        map<int, int[]> var_to_clause
        map<int, (int, int)> clause_to_vars
\end{lstlisting}

The \textit{var\_to\_trail} field is the only change we'll need to make, and it 
is just a mapping from variable index to trail index. If you enforce the invariant 
that each variable maps to exactly one place in the trail and no two variables 
map to the same place in the trail at the same time, you don't need to initialize 
these values to anything. 

The next thing we'll look at is our \textit{propagate()} code since it needs very 
minimal changes. The change we're making now is that instead of returning a boolean
value to indicate whether or not propagation was successful, we're now going to return 
a clause if we encounter a contradiction and nothing otherwise.

\begin{lstlisting}
    bool propagate(int[] assignment, int decisionLevel) {
        while (trailHead < trail.size) {
            int prop = trail[trailHead].lit
            trailHead = trailHead + 1
            ...
            for (i = 0 up to watched_clauses.size) {
                clause_idx = watched_clauses[i]
                int[] clause = clauses[clause_idx]
                (int, int) watches = clause_to_var[clause_idx]
                ...
                if (relocated == False) {
                    ...
                    int other_watch = clause[watches.second];
                    int idx = abs(other_watch)
                    if (assignment[idx] == Unassigned) {
                        ...
                        append(trail, {other_watch, decisionLevel, clause_idx})
                    } else {
                        return clause
                    }
                }
            }
        }
        return null
    }
\end{lstlisting}

And before we look at changes to our solving method, we're going to create a new 
overload for \textit{createWatched()} that takes in the index of a clause. Its 
only purpose is to be called on newly learned clauses, which don't have the same 
guarantees as the other overload. You see, the learned clause should hopefully 
NEVER be unsatisfied when this method operates over it. It should ALWAYS return 
True, so if it ever doesn't, it's a good debugging tool. Here's how it works:
we traverse over the clause of the passed in index, counting the number of 
True, False, and Unassigned literals it has. We will then search for two 
indices that could be watched literals. Since the clause should never be 
unsatisfied, we are guaranteed to always find one (if our implementation is correct).
If we can't find another and our clause has more than one literal, we will initialize 
the second watched literal to any arbitrary place in the clause (with the caveat 
that it can't be the same index as our first watched literal). We will then check if the 
clause is unit, and if it is, we will add it to our trail, assign the Unassigned variable,
etc. Otherwise, we do nothing.

\begin{lstlisting}
    bool createWatched(int index, int decisionLevel) {
        clause = clauses[index]
        if (clause.size == 0) {
            return False
        }

        bool foundFirst = False 
        bool foundSecond = False 
        int firstIdx = -1
        int secondIdx = -1

        int numUnassigned = 0
        int numFalse = 0
        int numTrue = 0

        for (i = 1 up to clause.size) {
            int idx = abs(clause[i])
            if (assignment[idx] == Unassigned) {
                numUnassigned = numUnassigned + 1
                if (foundFirst == False) {
                    foundFirst = True
                    firstIdx = i
                } else if (foundSecond == False) {
                    foundSecond = True 
                    secondIdx = i
                }
            } else if (evals_false(clause[i])) {
                numFalse = numFalse + 1
            } else {
                numTrue = numTrue + 1

                if (foundFirst == False) {
                    foundFirst = True
                    firstIdx = i
                } else if (foundSecond == False) {
                    foundSecond = True 
                    secondIdx = i
                }
            }
        }

        if (clause.size >= 2) {
            if (foundSecond == False) {
                secondIdx = (firstIdx + 1) % clause.size
            }

            append(var_to_clause[clause[firstIdx]], index)
            append(var_to_clause[clause[secondIdx]], index)
            
            clause_to_var[index] = {firstIdx, secondIdx}
        } else {
            append(var_to_clause[clause[firstIdx]])
            clause_to_var[index] = {firstIdx, firstIdx}
        }

        if (numTrue == 0 && numUnassigned == 1) {
            append(trail, {clause[firstIdx], decisionLevel, index})
            var_to_trail[abs(clause[firstIdx])] = trail.size() - 1

            int var = clause[firstIdx]
            int idx = abs(var)

            if (var > 0) {
                assignment[idx] = True
            } else {
                assignment[idx] = False
            }
        }

        return True
    }
\end{lstlisting}

Because we don't know anything about this clause, we need to gather as much 
information as we can right away before we decide how watched literals are 
spread out on the clause. When the clause has more than one literal, we 
use ``secondIdx = (firstIdx + 1) \% clause.size" because that prevents us 
from adding a watched literal past the end of the clause, if \textit{firstIdx}
is already on the edge.

The last method we need to discuss is our \textit{solve} loop. Clause 
resolution happens after backtracking because it relies on the current 
trail to understand
\footnote{It follows a mathematical process to transform the conflict clause 
into one that prevents that same conflict from ever occurring again. The 
algorithm does not ``understand'' anything. Do not anthropomorphize computer 
algorithms, especially artificial intelligence.}
why a conflict occurred. If we were to undo propagations, then the resolution 
operator would have no reference as to what ``most recently'' means, and 
therefore, clause resolution wouldn't work.

\begin{lstlisting}
    bool solve() {
        var[1..numVars] assignment

        for each (var in assignment) {
            var = Unassigned
        }

        if (createWatched(assignment) == False) {
            return False
        }

        int decisionLevel = 0

        if (trail.size > 0) {
            append(trailStarts, {0, true})

            if (propagate(assignment, decisionLevel) == False) {
                return False
            }
        }

        while (True) {
            if (trailHead == trail.size) {
                //Assignment code
                int firstUnassigned = -1
                for (i = 1 up to numVars) {
                    if (assignment[i] == Unassigned) {
                        firstUnassigned = i
                        break
                    }
                }

                if (firstUnassigned == -1) {
                    return True
                }

                assignment[firstUnassigned] = True
                append(trail, {firstUnassigned, -1})
                append(trailStarts, {trail.size - 1, False})
                trailHead = trailStarts.last.idx
            }

            optional int[] contradiction = propagate(assignment)

            if (contradiction has a value) {
                int[] conflict = contradiction.value

                bool propagated_found = True
                while (True) {
                    propagated_found = False
                    int idx = 0

                    for (i = 1 up to conflict.size) {
                        int var = conflict[i]
                        int v_idx = abs(var)

                        int t_idx = var_to_trail[v_idx]
                        if (trail[t_idx].reasonIdx has a value) {
                            propagated_found = True
                            idx = max(idx, t_idx)
                        }
                    }

                    if (propagated_found == False) {
                        break
                    }

                    int lit = trail[idx].lit
                    int[] reason = clauses[trail[idx].reasonIdx]

                    append(conflict, reason)

                    for (i = 0 up to conflict.size) {
                        if (abs(conflict[i]) == abs(lit)) {
                            swap(conflict[i], conflict.last)
                            erase(conflict, conflict.last)
                        } else {
                            i = i + 1
                        }
                    }

                    set<int> duplicates
                    for (i = 0 up to conflict.size) {
                        if (contains(duplicates, conflict[i])) {
                            swap(conflict[i], conflict.last)
                            erase(conflict, conflict.last)
                        } else {
                            insert(duplicates, conflict[i])
                            i = i + 1
                        }
                    }
                }

                while (trailStarts.size > 0) {
                    DecisionLevel decision = trailStarts.last

                    while (trail.size > decision.idx + 1) {
                        int lit = trail.last.lit
                        assignment[abs(lit)] = Unassigned
                        var_to_trail[abs(lit)] = -1
                        erase(trail, trail.last)
                    }

                    if (decision.triedFalse == False) {
                        TrailEntry entry = trail.last

                        entry.lit = -entry.lit
                        assignment[abs(entry.lit)] = False
                        var_to_trail[abs(entry.lit)] = -1
                        decision.triedFalse = True

                        trailHead = decision.idx
                        break
                    } else {
                        int lit = trail.last.lit
                        assignment[abs(lit)] = Unassigned
                        var_to_trail[abs(lit)] = -1
                        erase(trail.last)
                        erase(trailStarts.last)
                    }
                }

                if (trailStarts.size == 0) {
                    return False
                }

                append(clauses, conflict)
                createWatched(clauses.size)

                continue
            }
        }
    }
\end{lstlisting}

I want to point out that this isn't nevessarily CDCL yet. This is more like DPLL, but with 
adding clauses to our database. The reason as to why this is, is because whenever we recover 
from a contradiction, we only go back one level. You see, CDCL does away with the idea of a
``tried other branch" flag and relies on the learned clauses to guide what variable assignments
are possible and which ones are contradictory. However, because we only backtrack one level, there 
is nothing to actually force new variable assignments in cases where the learned clause isn't a unit 
clause. It can happen, but it doesn't always happen, so we still need to try the other assignment 
since for the most part, the learned clause isn't used.